%%%%%%%%%%%%%%%%%%%%%%%%%%%%%%%%%%%%%%%%%
% Medium Length Professional CV
% LaTeX Template
% Version 2.0 (8/5/13)
%
% Content adapted from main.tex while retaining the cv_4.tex layout.
%%%%%%%%%%%%%%%%%%%%%%%%%%%%%%%%%%%%%%%%%

\documentclass{resume}

\usepackage[dvipsnames]{xcolor}
\definecolor{TitleColor}{HTML}{17324D}
\usepackage{fontawesome5}
\usepackage{enumitem}
\usepackage{mathpazo}
\usepackage{tabularx}
\usepackage{titlesec}
\usepackage{changepage}
\usepackage{needspace}
\usepackage[
  colorlinks=true,
  urlcolor=TitleColor,
  linkcolor=TitleColor,
  pdfauthor={Shufan Shahi},
  pdftitle={Shufan Shahi - Curriculum Vitae}
]{hyperref}

\usepackage[left=0.5in,top=0.4in,right=0.5in,bottom=0.5in]{geometry}

\setlist[itemize]{
  leftmargin=1.25em,
  itemsep=1pt,
  topsep=2pt,
  parsep=0pt,
  partopsep=0pt
}
\setlength{\parskip}{2pt}
\setlength{\parindent}{0pt}

\newcommand{\tab}[1]{\hspace{.2667\textwidth}\rlap{#1}}
\newcommand{\itab}[1]{\hspace{0em}\rlap{#1}}
\newcommand{\projectheading}[4]{%
  \Needspace{3\baselineskip}% Keep the heading with at least the first detail line.
  \begin{tabularx}{\linewidth}{@{}X@{\hspace{1em}}r@{}}
    \textbf{#1 - #2} $|$ \emph{#3} & {\small #4}
  \end{tabularx}\vspace{-4pt}
}

\name{Shufan Shahi}

% resume.cls prints the third \address first, followed by the first and second.
\address{
  \faPhone\ \href{tel:+8801990935044}{+880 1990-935044} \quad
  \faEnvelope\ \href{mailto:shufanshahi@gmail.com}{shufanshahi@gmail.com}
}
\address{
  \faGithub\ \href{https://github.com/shufanshahi}{github.com/shufanshahi} \quad
  \faLinkedin\ \href{https://www.linkedin.com/in/shahi-shufan-680135289/}{linkedin.com/in/shahi-shufan} \quad
  \faGlobe\ \href{https://shufanshahi.github.io/}{shufanshahi.github.io}
}
\address{\faMapMarker*\ Dhaka, Bangladesh}

% Keep the template's centered header while containing its harmless \centerline warning.
\renewcommand{\printaddress}[1]{%
  \begingroup
    \hfuzz=31pt
    \def\\{\addressSep\ }%
    \centerline{#1}\par
  \endgroup
  \addressskip
}

\makeatletter
\renewcommand{\printname}{%
  \begingroup
    \hfuzz=31pt
    \hfil{\color{TitleColor}\MakeUppercase{\namesize\bfseries\@name}}\hfil
    \nameskip\break
  \endgroup
}
\makeatother

\titleformat{\section}
  {\Large\bfseries\scshape}{}{0em}{\color{TitleColor}}
  [\color{TitleColor}\titlerule]
\titlespacing*{\section}{0pt}{1.1em}{0.65em}

\renewenvironment{rSection}[1]{
  \Needspace{3\baselineskip}% Avoid an orphaned section title without reserving a large block.
  \section{#1}
  \begin{list}{}{
      \setlength{\leftmargin}{0.2em}
  }
  \item[]
}{
  \end{list}
  \vspace{0pt}
}

\begin{document}
\fontsize{10.5pt}{12.6pt}\selectfont

% \begin{rSection}{Professional Profile}
% Computer Science undergraduate with hands-on experience in data analysis, machine learning, and full-stack development. Built web and mobile applications using React, Node.js, and Firebase. Skilled in relational databases, API integration, and delivering real-world solutions through hackathons and academic projects.
% \end{rSection}

\begin{rSection}{Education}
\textbf{Islamic University of Technology} \hfill \textbf{2022--2026} \\
B.Sc. in Computer Science and Engineering \hfill \textbf{CGPA: 3.58} \\
Coursework: Algorithms, Machine Learning, Deep Learning, Databases, and Networks.

\textbf{Chuadanga Govt. College} \hfill \textbf{2021} \\
Higher Secondary Certificate (HSC) \hfill \textbf{GPA: 5.00/5.00}
\end{rSection}

\begin{rSection}{Experience}
\textbf{AI Engineer} \hfill \textbf{2026--Present} \\
\emph{GLiTch LABS}
\begin{itemize}
  \item Built RAG architectures, connecting LLMs to vector databases for optimized semantic search.
\end{itemize}

\textbf{Student Researcher} \hfill \textbf{2026--Present} \\
\emph{Elite Research Lab}
\begin{itemize}
  \item Support research projects through literature reviews, experiments, and analysis.
\end{itemize}

\textbf{Researcher} \hfill \textbf{2025--2026} \\
\emph{IUT Computer Vision Lab}
\begin{itemize}
  \item Researched multimodal AI methods to improve integration and mitigate modality imbalance.
\end{itemize}


\textbf{Senior Executive} \hfill \textbf{2025--2026} \\
\emph{Project Aero}
\begin{itemize}
  \item Contributed to drone R\&D by applying AI-based object tracking for international UAV competitions.
\end{itemize}



\end{rSection}



\begin{rSection}{Research Experience}
\Needspace{3\baselineskip}
\textbf{EthMuSAM : Mutual Scoring With SAM For Zero Shot Anomaly Segmentation} \hfill \textbf{2026}
\begin{itemize}
  \item Extracted patch-level representations from images using pretrained DINOv2 encoders.
  \item Localized anomalous regions by comparing test-image embeddings with representations from normal reference images.
\end{itemize}
\Needspace{3\baselineskip}
\textbf{Prototype-Guided Adaptive Feature Learning for Multimodal Emotion Recognition in Conversation (MERC)} \hfill \textbf{2026}
\begin{itemize}
  \item Developed graph-compatible adaptive feature weighting and modality-specific prototype supervision to learn class-discriminative representations before cross-modal reasoning.
  \item Improved multimodal emotion-recognition performance across multiple benchmark datasets through adaptive feature processing and prototype-guided learning.
\end{itemize}
\Needspace{3\baselineskip}
\textbf{Rethinking Visual Features in Graph Based Multimodal Emotion Recognition in Conversations (MERC): From Low Level Pixels to Semantic Representations} \hfill \textbf{2026}
\begin{itemize}
  \item Developed a semantic visual-feature pipeline that converts utterance-level videos into temporally structured descriptions of facial expressions, gaze, body language, and emotional transitions.
  \item Encoded the descriptions into dense representations for graph-based multimodal reasoning, improving cross-modal alignment while providing a viable alternative to noisy handcrafted visual features.
\end{itemize}

\Needspace{3\baselineskip}
\textbf{Investigating Lightweight Attention Placement in MobileNetV3 for Image Classification on CIFAR-100} \hfill \textbf{2026}
\begin{itemize}
  \item Ran a controlled six-variant ablation of ECA attention in MobileNetV3-Small, comparing addition, stacking, and replacement of native SE modules with paired significance testing across seeds.
  \item Replacing SE with ECA cut parameters by 28.13\% (1.61M to 1.16M) with comparable accuracy (70.62\% vs. 70.77\% baseline).
\end{itemize}
\Needspace{3\baselineskip}
\textbf{Deepfake Detection Using Cross-Covariance Image Transformers (XCiT)} \hfill \textbf{2026}
\begin{itemize}
  \item Trained an XCiT-based classifier to distinguish real from synthetic face images on a DeepfakeBench-derived dataset of over 260K images.
  \item Achieved 94.4\% validation accuracy using linear-complexity cross-covariance attention.
\end{itemize}

\end{rSection}

\Needspace{6\baselineskip}
\begin{rSection}{Projects}


\projectheading{NeuroEye}{Retinal Disease Classification}{PyTorch, EfficientNetV2-L, ResNet-50, ViT}{}
\begin{itemize}
  \item Built a weighted ensemble of EfficientNetV2-L, ResNet-50, and ViT-B/16 with learnable fusion weights to classify 10 retinal conditions from fundus images.
  \item Trained with mixed precision, per-model learning rates, and early stopping; reached 77.9\% test accuracy and 0.80 macro F1.
\end{itemize}

\projectheading{Drone Vision}{Aerial Tracking}{PyTorch, YOLOv11X, ByteTrack}{}
\begin{itemize}
  \item Fine-tuned YOLOv11X on VisDrone and integrated ByteTrack for human counting and multi-object tracking in aerial footage.
\end{itemize}

\projectheading{Benhallu}{Bengali Hallucination Detection}{Python, RAG, BanglaBERT, scikit-learn}{}
\begin{itemize}
  \item Built a retrieval-augmented pipeline to verify Bengali answers against supplied context or a Bengali QA knowledge base.
  \item Created a hallucination dataset and used fine-tuned BanglaBERT with answer-comparison features to classify responses.
\end{itemize}

\projectheading{Bengali Speaker Diarization}{Speaker Turn Detection}{PyTorch, Silero VAD, ECAPA-TDNN}{}
\begin{itemize}
  \item Built a speech pipeline using neural voice activity detection, speaker embeddings, and agglomerative clustering to label timestamped speaker turns.
  \item Estimated speaker counts with silhouette scores and exported diarization results to CSV.
\end{itemize}





\end{rSection}

\begin{rSection}{Technical Skills}
\textbf{Programming Languages:} Python, C++, Java \\
\textbf{Data and Machine Learning:} PyTorch, TensorFlow, scikit-learn, NumPy, Pandas \\
\textbf{Databases:} MySQL, Firebase, PostgreSQL \\
\textbf{Web Development:} React, Tailwind CSS, Next.js, Spring Boot \\
\textbf{Mobile Development:} React Native, Flutter \\
\textbf{Visualization:} Matplotlib, Seaborn \\
\textbf{Development Tools:} Git, Linux \\
\end{rSection}


\begin{rSection}{Achievements, Awards and Scholarships}
\begin{itemize}[leftmargin=*,label=$\bullet$]
    \item \textbf{5th Place, VAND4 International Competition} \hfill \textbf{2026} \\
  Zero-Shot Industrial Track; built a zero-shot computer vision pipeline for detecting anomalies in industrial images using foundation models.
  
  \item \textbf{First Runner-up, IUT Coderush Hackathon} \hfill \textbf{2025} \\
  Built an application under tight deadlines in a fast-paced team.
\item \textbf{7th Place, BAIC Datathon} \hfill \textbf{2025} \\
  Analyzed datasets using preprocessing, visualization, and modeling.
\item \textbf{IUT Partial Scholarship} \hfill \textbf{2022} \\
  Secured a partial scholarship covering all four years of B.Sc. studies.
\item \textbf{Runner-up, State-Level Physics Olympiad} \hfill \textbf{2018}
\end{itemize}
\end{rSection}



\begin{rSection}{Reference}
\textbf{Shahriar Ivan} \\
Assistant Professor, Department of CSE, IUT \\
\textbf{Specialization:} Machine Learning, Deep Learning, Computer Vision \\
\faEnvelope\ \href{mailto:shahriarivan@iut-dhaka.edu}{shahriarivan@iut-dhaka.edu}
\end{rSection}

\end{document}
